\documentclass[11pt]{article}
\usepackage[margin=0.85in]{geometry}
\usepackage{graphicx}
\usepackage{booktabs}
\usepackage{array}
\usepackage{amsmath}
\usepackage{xcolor}
\usepackage[hidelinks]{hyperref}
\usepackage[numbers]{natbib}
\usepackage{microtype}

\definecolor{accent}{HTML}{2F6F7E}
\hypersetup{colorlinks=true,linkcolor=accent,citecolor=accent,urlcolor=accent}
\setlength{\parindent}{0pt}
\setlength{\parskip}{0.55em}

\title{\textbf{8\_16\_2026\_Update}\\
Binding Affinity Predictor}
\author{yqc5415}
\date{August 2026}

\begin{document}
\maketitle

\section*{Executive summary}

The long-term objective is a controlled comparison of Boltz-2, Nesso-1, and
FlashBind: what each model receives, what quantity it predicts, and how its
claims survive a common benchmark. Boltz-2 motivates particular care because
cofolded structure confidence and its affinity-related outputs are distinct
objects \citep{passaro2025boltz2}.

This update first develops a high-level picture of Boltz-2 and then uses it to
explain Nesso-1. Boltz-2 combines a learned representation trunk, full-atom
diffusion, confidence estimation, and affinity prediction. Nesso-1 retains a
coarse representation of protein--ligand geometry while removing the expensive
full-atom decoder, making it primarily a rapid affinity-screening model
\citep{shenoy2026nesso1}.

\section{Scientific questions}

\begin{enumerate}
  \item What molecular information do Boltz-2 and Nesso-1 receive?
  \item How does each model represent the protein, ligand, and their possible
        geometry?
  \item Why does Boltz-2 generate full-atom coordinates while Nesso-1 does not?
  \item What do their binder and continuous-affinity outputs mean?
  \item What physical detail does Nesso-1 sacrifice to obtain speed?
\end{enumerate}

\section{Boltz-2 architecture: a high-level view}

Boltz-2 combines two related tasks: it first predicts a full-atom
biomolecular complex and then uses the predicted interface to estimate whether
the ligand binds and how strong that binding may be \citep{passaro2025boltz2}.
The model is easiest to understand as four connected components: a trunk, a
diffusion decoder, a confidence module, and an affinity module.

\begin{center}
\fbox{\parbox{0.92\textwidth}{\centering
Molecular inputs
$\longrightarrow$ representation trunk
$\longrightarrow$ full-atom diffusion
$\longrightarrow$ candidate 3D complexes
$\longrightarrow$ confidence ranking
$\longrightarrow$ binder probability + continuous affinity score}}
\end{center}

\subsection{Inputs and molecular representation}

For a protein--ligand example, the basic inputs are the protein sequence and a
ligand description such as SMILES. The model can also use a protein
multiple-sequence alignment, structural templates, experimental-method labels,
and user-provided contact or pocket constraints. An atom-attention encoder
combines detailed atomic features into token representations. Protein residues
are represented as residue-level tokens, whereas ligand heavy atoms are
represented individually.

The trunk maintains two complementary learned objects:

\begin{itemize}
  \item a \emph{single representation} $s_i$ describing each token; and
  \item a \emph{pair representation} $z_{ij}$ describing the relationship
        between every pair of tokens.
\end{itemize}

Template and MSA information are incorporated before a 64-layer PairFormer
updates these representations. Recycling sends the updated representations
through the trunk repeatedly so that sequence, molecular identity, and
putative geometry can be refined together.

\subsection{Full-atom diffusion decoder}

The trunk does not itself provide the final coordinates. Boltz-2 begins the
structure-generation stage with noisy, effectively random atomic coordinates
and applies a learned denoising module repeatedly. At a schematic level,
\[
  \mathbf{x}_{t}
  \longrightarrow
  \mathbf{x}_{t-1}
  \longrightarrow \cdots \longrightarrow
  \mathbf{x}_{0},
\]
where the trunk representations condition every denoising step. The final
$\mathbf{x}_0$ contains explicit coordinates for the modeled atoms in the
protein--ligand complex. Because the initial coordinates and sampling process
are stochastic, several candidate complexes can be generated from the same
input.

Optional steering potentials can bias this reverse-diffusion process toward
user constraints or improved physical plausibility. These potentials guide
sampling; they do not turn the neural decoder into a molecular-dynamics or
free-energy calculation.

\subsection{Confidence module}

A separate, smaller PairFormer examines the final trunk representation and a
predicted coordinate set. It estimates confidence quantities such as local and
pairwise errors and interface confidence. These values help rank alternative
structure samples. Confidence answers ``how trustworthy is this predicted
geometry?''; it is not itself a measurement of binding strength.

\subsection{Affinity module}

The affinity calculation uses both the learned trunk representation and the
predicted structure. In the reported inference protocol, five structures are
generated and the candidate with the highest protein--ligand interface
confidence is selected. A specialized PairFormer then focuses on
protein--ligand and intra-ligand pairs, combines $z_{ij}$ with a distogram
derived from the selected coordinates, and pools the interface into one global
representation.

Two multilayer-perceptron heads produce distinct outputs:

\begin{enumerate}
  \item a binding-likelihood probability intended for binder/non-binder
        classification; and
  \item a continuous, logarithmic affinity value intended mainly for ranking
        related ligands against the same target.
\end{enumerate}

The continuous head is trained on a mixture of $K_i$, $K_d$, and IC$_{50}$
measurements standardized to micromolar units. It should be interpreted as an
IC$_{50}$-like measure of binding strength, not as an equilibrium constant from
one uniform assay and not as a calculated $\Delta G_{\mathrm{bind}}$.

\subsection{What Boltz-2 gives us---and what it does not}

Boltz-2 therefore provides two kinds of predictions that must remain separate:

\begin{center}
\begin{tabular}{p{0.27\textwidth}p{0.61\textwidth}}
\toprule
\textbf{Output family} & \textbf{Scientific interpretation} \\
\midrule
Structure and confidence & A sampled full-atom complex and the model's
estimated uncertainty in that geometry \\
Binder likelihood & A learned classification probability for interaction \\
Continuous affinity & A heterogeneous-assay, IC$_{50}$-like ranking score \\
\bottomrule
\end{tabular}
\end{center}

The affinity estimate can fail when the model selects the wrong pocket,
misplaces the ligand or protein conformation, or omits important context. The
released affinity module does not explicitly model contributions from water,
ions, essential cofactors, or all binding partners. Full-atom coordinates make
mechanistic inspection possible, but they do not guarantee physical validity
or thermodynamic accuracy.


\section{Nesso-1 architecture: a high-level view}

Nesso-1 asks whether binding affinity can be estimated from an approximate map
of the protein--ligand interface without first generating every atomic
coordinate. It is therefore best understood as a \emph{coarse-grained
cofolding affinity model}, rather than as a full-atom structure predictor
\citep{shenoy2026nesso1}.

\begin{center}
\fbox{\parbox{0.92\textwidth}{\centering
Protein sequence + ligand SMILES
$\longrightarrow$ protein and ligand embeddings
$\longrightarrow$ pairwise distance representation
$\longrightarrow$ likely binding-region crop
$\longrightarrow$ binder probability + continuous affinity score}}
\end{center}

\subsection{Inputs and molecular tokens}

The inputs are a protein amino-acid sequence and a ligand SMILES string. No
experimental complex, user-specified pocket, or multiple-sequence alignment is
required. The protein receives contextual embeddings from the pretrained
ESM-2 protein language model. Each protein residue is represented
geometrically by one center, its C$\beta$ atom (C$\alpha$ for glycine), while
each ligand heavy atom remains an individual token. Thus, the protein is
strongly coarse-grained but the ligand retains its heavy-atom graph.

\subsection{A learned map of pairwise geometry}

The central object is a pair representation $z$. For every pair of tokens,
$z_{ij}$ stores a learned description of their relationship. Protein--ligand
entries answer questions such as ``is residue 42 likely to be close to ligand
atom 5?'' A streamlined PairFormer updates this matrix using triangle
attention and triangle multiplication, operations that allow relationships
among three tokens to constrain one another.

The trunk predicts a categorical distance distribution,
\[
  P(d_{ij}\in\text{distance bin }k),
\]
rather than one exact coordinate set. From this distribution, Nesso-1 obtains
an expected distance and an entropy that measures uncertainty in the predicted
protein--ligand geometry.

\subsection{Finding the relevant protein region}

The first trunk pass considers the full protein. Subsequent recycling steps
retain protein tokens predicted to lie within approximately 22~\AA{} of the
ligand. The affinity module uses a tighter crop: residues predicted to lie
within approximately 15~\AA{} of any ligand heavy atom. In this sense, Nesso-1
infers a likely interaction region internally rather than requiring the user to
provide a pocket.

\subsection{Affinity outputs}

A second, smaller PairFormer---the affinity module---combines the protein
language-model embeddings, ligand features, pair representation $z$, and its
predicted distance information. It produces two distinct outputs:

\begin{enumerate}
  \item a binder likelihood for hit-identification classification; and
  \item a continuous affinity or potency score for compound ranking.
\end{enumerate}

The continuous output is not a calculated $\Delta G_{\mathrm{bind}}$. During
training, $K_i$, $K_d$, IC$_{50}$, and EC$_{50}$ labels are treated
interchangeably to accommodate incomplete and heterogeneous public assay
metadata. It should therefore be interpreted as an assay-trained
affinity-like score, with particular value for ranking compounds measured in
comparable settings.

\subsection{The essential contrast with Boltz-2}

\begin{center}
\begin{tabular}{p{0.21\textwidth}p{0.34\textwidth}p{0.34\textwidth}}
\toprule
 & \textbf{Boltz-2} & \textbf{Nesso-1} \\
\midrule
Structural decoder & Full-atom diffusion producing explicit coordinates & No
heavy-atom diffusion; predicts coarse pair distances \\
Protein context & MSA can be used & ESM-2 embeddings; no MSA \\
Primary structural output & Full biomolecular complex & Coarse interface
representation; approximate poses can be reconstructed \\
Design priority & Detailed structure plus affinity & Rapid affinity screening \\
\bottomrule
\end{tabular}
\end{center}

Removing full-atom diffusion is the main source of Nesso-1's speed. The cost is
that it cannot resolve side-chain conformations, hydrogen-bond geometry,
steric clashes, or water-mediated interactions with atomistic precision. Its
architectural hypothesis is that coarse interface geometry may be sufficient
for rapid affinity ranking, even though full physical detail remains important
for mechanistic interpretation and late-stage optimization.


\section{Next steps}

\begin{enumerate}
  \item Develop an equally clear architecture summary for FlashBind.
  \item Define a common native-output schema without silently equating binder
        probability, heterogeneous potency labels, and binding free energy.
  \item Choose a small benchmark with exact protein constructs, ligand
        stereochemistry, assay identity, qualifiers, and experimental units.
\end{enumerate}

\section*{Reproducibility}

The architecture descriptions are grounded in the downloaded version-1
technical reports for Boltz-2 and Nesso-1. The accompanying manifest records
the project Git revision, Overleaf revision, citation keys, and compiled PDF
checksum.

\bibliographystyle{unsrtnat}
\bibliography{references}

\end{document}
